\documentstyle[12pt,fullpage,steve]{article}
\setcounter{page}{1}

\begin{document}

\title{EE 150 Final Exams : Part II \\ {\bf CLOSED BOOK}}
\author{Wednesday, December 11 1996}
\date{}
\maketitle

\begin{itemize}
\item {\bf Name:}
\item {\bf Student Id Number:}
\end{itemize} 

{\bf Read First:}
\begin{enumerate}
\item [A.] {\bf **READ** THE DIRECTIONS CAREFULLY!}
\item [B.] Any materials including your text and notes may be consulted
during the exam at your discretion. 
\item [C.] You have a maximum of 50 minutes to complete the exam.  The exam
paper will be picked up from your desk (do not get up) at the end of
the alloted time.
\item [D.] There are 3 problems worth a total of 50 points. Be careful not to
spend too much time on any particular problem.  {\bf If you get stuck,
go on to the next problem}.  You can always come back to the ones that
are giving you problems.   
\item [E.] Write all answers, including any scratch paper or rough 
notes on the pages given you.
\item [F.] If you split the exam paper, write your student id number
on each and every page.
\end{enumerate}

% ---------------------------------------------------------------------------
\pagebreak
{\bf PROBLEM \#1} (20 points)

Below is a C function named {\em replace}.  {\bf It doesn't work
like it supposed to}.

\begin{verbatim}
int replace(int x[], int old, int new, int stop)  {
  int i;
  printf("Replace Call: old = %i, new = %i, stop = %i\n", old, new, stop);*/
  for (i = 1; x[i] >= stop; i++)
    {
      printf("Replace Loop: i = %i, x[i] = %i\n", i, x[i]);*/
      if (x[i] != old)
        {
          x[i] = new;
          printf("Replace New : i = %i, x[i] = %i\n", i, x[i]);
        }
    }
  printf("Replace Return value is i = %i\n", i);
  return i;
}
\end{verbatim}

What it is supposed to do is this: run through the
first $stop$ elements of array $x$, replacing each element that
contains the value $old$ with the value $new$.  It is then supposed to
return the number of elements it modified.  So, for example, suppose 
$table$ was array of $int$s that was initialized
this way:
\begin{verbatim}
int table[] =
  {31, 17, 26, 9, 6, 17, 19, 17, 21, 17, 96, 108, 61, 19, 17};
\end{verbatim}
Then the code:
\begin{verbatim}
n = replace(table, 17, 40, 9);
printf("%i\n", n);
\end{verbatim}
should have written a 3 and changed table to look like:
\begin{verbatim}
   31, 40, 26, 9, 6, 40, 19, 17, 21, 17, 96, 108, 61, 19, 17
\end{verbatim}
That is, it should have run through the first 9 elements of <TT>table</TT>,
changing each 17 to a 40.  But it doesn't!  

{\bf continued on next page ...}

% ---------------------------------------------------------------------------
\pagebreak
{\bf PROBLEM \#1-A} (7 points)

Trace through the program shown using the Tables model shown in class
and used repeatedly on quizzes.  If you do not use the Tables model
{\bf you will receive 0 on this part.}

Write your solution below.  Explicitly state any assumptions you make.

% ---------------------------------------------------------------------------
\pagebreak
{\bf PROBLEM \#1-B} (3 points)

What does this code actually write for the value of $n$ ?
\vspace{1.0in}

{\bf PROBLEM \#1-C} (5 points)

What does $table$ look like when the function is completed ?
\vspace{2.0in}

{\bf PROBLEM \#1-D} (5 points)

What has to be done to repair $replace$?  Either show a fixed version
of  the function, or explain {\em exactly} what must be changed and 
what the changes are.

\vspace{4.0in}

% ---------------------------------------------------------------------------
\pagebreak
{\bf PROBLEM \#2} (15 points)

Using the TABLES model, show how this program executes {\em and} what 
it prints out.

\begin{verbatim}
#include <stdio.h>
main()
{
int a, b;
int c; float d, f;
int *p1, *p2;  
int e = 10; 
double t[3] = {2.0,2.1,2.2};
double *p3;
double dd = 1.5;
int fizbin(int *, int *);
b = 2, a = b++;
d = 2.5, c = ++d;
dd = (double) (6.55 + 4.44) * (int) 3.1;
p1 = &e; p2 = p1, *p2 += 3;  *p1 = 11; a = *p1, *p2;  
*(&a) = *t, p3 = t + 2, d = a + *p3;
*(&(*(&d))) = *(t + 2), *(t + 2) = 3; f = t[2];
a = 5; b = 3; c = 1; c = (a == b) + (a > b) * (c < a) - (c < a > b);
for(c=1;c<3;c++) c = fizbin(&a,&b);
}

int fizbin(int *i,int *j) {
     *i = *i * *j;
     *j = *i - *j;
     return *i - *j; }
\end{verbatim}

Write your solution below.  Explicitly state any assumptions you make.

% ---------------------------------------------------------------------------
\pagebreak
{\bf PROBLEM \#2} ( continued ... )

% ---------------------------------------------------------------------------
\pagebreak
{\bf PROBLEM \#3} (15 points)

We are given the task of producing a histogram for a set of student
grades.  Imagine we are given the four cut off points for each of the
letter grades of A, B, C, and D.  In addition, we are given a number
indicating how many buckets (rows) should appear in a histogram, a
number indicating the maximum possible student score (the lowest is
0),  a number indicating how many student grades are following, and
then the list of student grades.

Each bucket covers an equal range of student grades.  For example, a
number of buckets equal 4 and maximum grade equal 100 would indicate
each bucket covers a range of 25 marks.

We would like to produce a histogram output showing the student score
together with the grade letter for each grade input. 

For example, an input file that looks like the following:
\begin{verbatim}
85 74 63 52 3 100 9 67 56 91 94 96 23 88 55 66
\end{verbatim}

Says we want 3 rows of output. Also, since the maximum score is 100, we know
the range of each column is 0-33, 33-66 and 66-100. The output would look like:
\begin{verbatim} 
0-33    23-F
33-66   56-D 55-D 66-C
66-100  67-C 91-A 94-A 96-A 88-A 
\end{verbatim}

Write a program that creates this histogram given the specified input.
Explicitly state any assumptions you make.

% ---------------------------------------------------------------------------
\pagebreak
{\bf PROBLEM \#3} ( continued ... )


\end{document}

/* Local Variables: */
/* compile-command: "latex ee150-final2" */
/* End: */